% Français 9115, batch 6 — Gallica views 101–120.
% Pass: Opus 5 (claude-opus-5), 2026-09-15 — first pass, unchecked against
% the views by a human.
%
% What the twenty views hold. Two things, and the seam falls at view 115.
% Views 101 to 114 carry on the working drafts that batch 5 began at view 93:
% single leaves in her fast, heavily corrected hand, each written on one side
% only, on the introduction to the memoir on the curvature of surfaces — the
% passage where the elastic solid is reduced to its surface. Views 115 to 120
% are something else entirely: three printed leaves, each pasted on a mount,
% pp. 81–86 of Legendre's « Mémoire sur la détermination des fonctions Y et Z
% qui satisfont à l'équation 4(x^n − 1) = (x − 1)(Y^2 ± nZ^2), n étant un
% nombre premier 4i ∓ 1 », « Lu à l'Académie, le 11 octobre 1830 », with the
% volume mark « T. XI. » at the foot of the first page. The text breaks off at
% the foot of p. 86, inside Exemple I (n = 31); view 121, outside this batch,
% goes on. Twelve views of the twenty carry text: 101, 103, 105, 107, 109,
% 112, 114, 115–120. Views 102, 104, 106, 108 and 110 are blank versos showing
% the writing of the recto through the paper; views 111 and 113 carry nothing
% but an ink number (the text of those two leaves is on the other side, views
% 112 and 114). None of these is transcribed beyond the numbers they carry.
%
% The drafts, leaf by leaf.
%
%   v. 101  Headed « Le solide doué des forces d'élasticité est un système de
%           points matériels mobiles », with « N° 1 » in the margin and an
%           isolated « A » (batch 5's view 95 has « Commencement de la p. A »).
%           The displacement of the material points is not only the change of
%           their mutual distance but the common motion of a point and its
%           neighbours; the restoring forces are proportional to the absolute
%           displacement, measured from fixed points.
%   v. 103  Marked « B ». A history of the question in the first person: Chladni's
%           experiments concerned only plates; the Académie's programmes used the
%           word « surface » without defining it; the memoir published in 1814
%           « par un des membres de l'accademie » treats only plane surfaces;
%           « le mémoire couronné qui fut présenté en 1815 » first attempts curved
%           surfaces, by the cylinder; « mes recherches sur la théorie des
%           surfaces publiées en 1821 » take the same example and generalise the
%           hypothesis proposed for plane surfaces; and she acknowledges how much
%           Savart's experiments helped her develop the consequences of her
%           principles. The dates are on the leaf.
%   v. 105  That the character of the « surface » is still not generally known;
%           a struck passage says that the objection that only a surface of
%           infinitely small thickness was in question would not have been put to
%           the author of the 1814 memoir, nor accepted by him, had this been
%           reflected on. Then what « surface » designates: a solid formed of
%           infinitely many layers; all the points on a normal move through equal
%           spaces; the figure is the same for all the layers and independent of
%           the thickness, as the identity of the nodal figures on the two faces
%           attests.
%   v. 107  Marked « D ». Her hypothesis concerned only the curvature of the
%           element of a plane surface, yet one operation of calculus drew from it
%           an equation « dont personne ne conteste l'exactitude »; were the solid
%           not so regarded, how could the equation of the surface become that of
%           the solid by multiplication by a coefficient depending on thickness?
%   v. 109  The motion perpendicular to the tangent planes is an express
%           reservation; thickness is treated apart, and variable thickness is
%           reduced to equal thickness; a struck sentence mentions a memoir long
%           since presented to the Académie. The spaces run through are measured
%           by the change of the distances to the tangent plane, and the mean
%           curvature « est ici de grand secours ».
%   v. 112  « Le cas dont je me suis spécialement occupée »: a thin solid of any
%           curvature, points moving perpendicular to the tangent planes, every
%           layer similar to the faces; that case is that of vibrating elastic
%           surfaces; the identity of the nodal figures on opposite faces follows
%           from the identity of their curvature, confirmed by Savart's
%           experiments, and is the distinctive sign of this kind of vibration.
%   v. 114  In a finer pen and paler ink, three equations for the curvatures
%           1/R, 1/R', 1/R'' of normal sections making angles A − a, A, A + a
%           with the section of greatest curvature, in terms of 1/r, 1/r', sin a
%           and cos 2A; the right side of (3) vanishes for a sphere, for a = 0 or
%           180, or for A = 45; and the rate at which curvature changes between
%           neighbouring sections, symmetric about the section of mean curvature.
%
% The printed piece is transcribed in full as printed, with its running heads,
% page numbers and footnote given in notes; its misprints (a missing subscript
% at view 119, √n where the rule for 4i − 1 asks √−n, also view 119) stand.
% Why a memoir of Legendre read in 1830 sits among these leaves is not said on
% them.
%
% The numbers on the leaves. All in ink, all on the side of the leaf that
% carries a number; none inferred.
%   The series of batch 5 runs on: 50 (v. 101), 51 (v. 103), 52 (v. 105),
%   53 (v. 107), 54 (v. 109), 55 (v. 111), « 55 bis » with a struck « 66 »
%   below it (v. 113), 56 (v. 115), 57 (v. 117), 58 (v. 119). The numbers of
%   the printed leaves are written on the mounted printed pages.
%   Letters beside the drafts, apparently her own sequence marks: « A » (v. 101),
%   « B » (v. 103), « D » (v. 107); none was found on v. 105 or v. 109.
%   Printed pagination: 82, 83, 84, 85, 86 on views 116–120; view 115, the
%   first page, is unnumbered; « T. XI. » and the signature mark « 11 » at the
%   foot of view 115, « 11. » at the foot of view 117.
% No pencilled foliation of the Bibliothèque was found on any view, so no
% \folio{} is written, as in batches 1 to 5.
%
% Notation. At view 114 the equations are set as written, with the primes as
% far as they can be read; the sign between 1/r and 1/r' in (3) is under an
% ink blot and left blank, and the second denominator of (1) is read only as R.
% In the printed piece Legendre's radical sign is set as \sqrt, his symbol
% (k/n) as a parenthesised fraction, his « . » for multiplication kept.
%
% Her hand. The drafts are the hand of batch 5's drafts, worse: overwritten,
% struck line by line and crossed out whole, with additions between the lines
% and in the left margin, some in paler ink or pencil. Where the final state of
% a sentence is not coherent (views 101, 105, 109) the words are given in the
% order of the lines and the note says what is written where. \marginal{}
% holds her marginal additions; her underlining is \emph{}.
%
% What could not be read. Fifteen places, all in the drafts, none guessed:
% two struck interlinear additions at view 101; a struck word in « les ___
% variété » and a word struck after « Savart » at view 103; the word after
% « Je crois utile d' » and two struck words in the passage on the layers at
% view 105; two words in the struck sentence on the memoir presented to the
% Académie and a struck word before « des distances » at view 109; the word
% before « choix », three struck words in the struck paragraph, and a struck
% word before « ingénieuses » at view 112. The printed pages were read whole.
%
% Nothing here was checked against any published transcription or edition.
\documentclass[11pt,a4paper]{article}
% The preamble shared by every transcription of the mathematicians' manuscripts in Gallica.
%
% It rests on one idea: **uncertainty compounds, it does not resolve.** A
% transcription that smooths over a doubtful word has destroyed the only thing
% it was made for — a reader must be able to tell, at every point, what was on
% the page and what was supplied. The apparatus macros below are therefore the
% critical apparatus, not decoration.
%
% The same commands are recognised by `scripts/render.mjs`, which produces the
% site's reading view, and by `scripts/tei.mjs`. A macro added here must be
% added there too, or rendering fails loudly — which is the wanted behaviour.
%
% Comments here are English, like the rest of the repository. The documents
% this preamble sets are French, and that is deliberate: see the skills.

\ProvidesPackage{germain}[2026/09/15 Transcriptions of mathematicians' manuscripts in Gallica]

\usepackage{amsmath,amssymb,amsthm}
\usepackage{mathrsfs}
% Kept from the parent preamble so the renderer's diagram support stays
% available; these manuscripts draw no commutative diagrams, and nothing here uses it.
\usepackage{tikz-cd}
\usepackage[english,french]{babel}
\usepackage{fontspec}

% --- Greek ------------------------------------------------------------------
% The text face stays fontspec's default, Latin Modern, which has no Greek:
% XeTeX drops every glyph it lacks with a line in the log and nothing on the
% page, and the Greek words the manuscripts quote vanished from the PDFs. So
% the Greek and Greek Extended blocks get a XeTeX character class of their own,
% and entering or leaving a run of it switches the family to CMU Serif — the
% same Computer Modern design, with polytonic Greek — and back. Only the
% family changes: a Greek word in \textit or \textbf keeps its shape and series.
% The transitions to and from every other class carry the tokens that class
% has towards an ordinary letter, so French spacing before « ; » or inside
% « » still holds after a Greek word. No grouping, so no brace or \endgroup
% can come out unbalanced; maths is left alone, since XeTeX inserts nothing
% there. Kept identical in `exercices.sty`.
\makeatletter
\newfontfamily\germainGreekfont{cmun}[
  Extension=.otf, NFSSFamily=germaingreek,
  UprightFont=*rm, ItalicFont=*ti, SlantedFont=*sl,
  BoldFont=*bx, BoldItalicFont=*bi, BoldSlantedFont=*bl]
\newXeTeXintercharclass\germain@greek
\def\germain@greekrange#1#2{%
  \@tempcnta=#1\relax
  \loop
    \XeTeXcharclass\@tempcnta=\germain@greek
    \ifnum\@tempcnta<#2\advance\@tempcnta\@ne\repeat}
\germain@greekrange{"0370}{"03FF}
\germain@greekrange{"1F00}{"1FFF}
\def\germain@greekprev{\rmdefault}
\def\germain@greekin{%
  \edef\germain@greekprev{\f@family}\fontfamily{germaingreek}\selectfont}
\def\germain@greekout{\fontfamily{\germain@greekprev}\selectfont}
\def\germain@greekjoin{%
  \XeTeXinterchartokenstate=\@ne
  \@tempcnta=\z@
  \loop
    \ifnum\@tempcnta=\germain@greek\else
      \XeTeXinterchartoks\germain@greek\@tempcnta=\expandafter{%
        \expandafter\germain@greekout\the\XeTeXinterchartoks\z@\@tempcnta}%
      \XeTeXinterchartoks\@tempcnta\germain@greek=\expandafter{%
        \the\XeTeXinterchartoks\@tempcnta\z@\germain@greekin}%
    \fi
    \ifnum\@tempcnta<4095\advance\@tempcnta\@ne\repeat}
% babel-french sets its own classes, and saves and restores the interchar
% state, each time a language is selected: join again after it does.
\addto\extrasfrench{\germain@greekjoin}
\addto\extrasenglish{\germain@greekjoin}
\AtBeginDocument{\germain@greekjoin}
\makeatother

\usepackage{geometry}
\geometry{a4paper,margin=25mm,marginparwidth=28mm}
\usepackage{marginnote}
\usepackage{xcolor}
\usepackage[normalem]{ulem}
\setlength{\emergencystretch}{3em}
\usepackage{hyperref}
\hypersetup{colorlinks=true,linkcolor=black,urlcolor=black,pdfborder={0 0 0}}

% --- Batch metadata ---------------------------------------------------------
% Taken as they stand by the rendering script, which shows them at the head of
% the reading view. They fix what the file claims to cover. The macro names are
% the parent project's, so that the scripts stay shared; \folder here means the
% volume's slug — `fr-9115` — and \pages counts Gallica views.

\newcommand{\folder}[1]{\gdef\grothFolder{#1}}
\newcommand{\batch}[1]{\gdef\grothBatch{#1}}
\newcommand{\pages}[2]{\gdef\grothFirst{#1}\gdef\grothLast{#2}}
\newcommand{\foldertitle}[1]{\gdef\grothFoldertitle{#1}}

% The holder's own shelfmark, printed as they write it: « Français 9115 ».
\newcommand{\shelfmark}[1]{\gdef\grothShelfmark{#1}}

% The dating, copied verbatim from the holder's catalogue — never inferred
% here. « XIXe siècle » is the BnF's; a pass that dates a leaf from its content
% says so in a \note{}, as its own inference.
\newcommand{\dating}[1]{\gdef\grothDating{#1}}

% The status notice, printed in the title block and repeated as a diagonal
% watermark on every page of the PDF. The manuscripts are in the public
% domain; what this line declares is not a right but a quality — an
% unverified first pass by a machine. The skills write the line; a file
% without it gets no watermark, so leaving it out is visible in the output.
\newcommand{\watermark}[1]{\gdef\grothWatermark{#1}}

\gdef\grothFolder{?}\gdef\grothBatch{}\gdef\grothFirst{?}\gdef\grothLast{?}%
\gdef\grothFoldertitle{}\gdef\grothShelfmark{}\gdef\grothDating{}\gdef\grothWatermark{}

% --- The résumé -------------------------------------------------------------
\newenvironment{resume}
  {\small\setlength{\parskip}{0.45em}}
  {\par\medskip\hrule\medskip}

% --- Keywords ---------------------------------------------------------------
% In English, closing the modernised reading's résumé: the modern vocabulary
% under which to search for what the leaves construct. The manifest extracts
% them to tag the volume — this line is the single source of the tags.
\newcommand{\keywords}[1]{\par\smallskip\noindent
  {\footnotesize\textsc{Keywords} --- \textit{#1}}\par}

% --- The critical apparatus -------------------------------------------------

% The Gallica view number, in the margin. This is the anchor that lets a
% disagreement over a formula be settled by looking at *one* image rather than
% a volume of 750. The site uses it to turn the facsimile as the transcription
% is scrolled. A view is one scanned image: usually one side of a leaf.
\newcommand{\page}[1]{%
  \marginnote{\footnotesize\color{gray!70}\textsf{v.\,#1}}%
  \label{page:#1}\ignorespaces}

% The BnF's pencilled foliation, where the leaf carries it and a pass can read
% it: « 198r ». This is what the literature cites — « ff. 198r–208v » — and
% the one line that turns a citation into a view. Recorded, never inferred:
% a folio a pass cannot read is not written.
\newcommand{\folio}[1]{%
  \marginnote{\footnotesize\color{gray!50}\textsf{f.\,#1}}\ignorespaces}

% The modernised reading's coarser counterpart to \page{}: which views a
% section is drawn from.
\newcommand{\pagerange}[2]{%
  \marginnote{\footnotesize\color{gray!70}\textsf{v.\,#1--#2}}%
  \ignorespaces}

% Illegible: never guessed.
\newcommand{\ill}{\textcolor{red!60!black}{[\,\ldots\,]}}

% A doubtful reading: the word is given, and flagged as uncertain.
\newcommand{\uncertain}[1]{\uline{#1}}

% An editorial addition — a missing letter, an implied word.
\newcommand{\add}[1]{\textcolor{blue!50!black}{[#1]}}

% Struck out by the writer. What was crossed out is part of the document.
\newcommand{\struck}[1]{\sout{#1}}

% The transcriber's note, on the state of the page or the mathematics.
\newcommand{\note}[1]{\par\smallskip{\small\color{gray!60!black}\textit{[#1]}}\par\smallskip}

% A marginal note of hers, distinct from the body of the text.
\newcommand{\marginal}[1]{\par\smallskip\noindent\hspace{1em}%
  {\small\color{gray!60!black}#1}\par\smallskip}

% --- Title ------------------------------------------------------------------
% The title block is French, like the documents themselves.

\AddToHook{shipout/background}{%
  \ifx\grothWatermark\empty\else
    \begin{tikzpicture}[remember picture, overlay]
      \node[rotate=52, scale=3.4, align=center, text=gray!18,
            font=\sffamily\bfseries]
        at (current page.center) {\grothWatermark};
    \end{tikzpicture}%
  \fi}

\AtBeginDocument{%
  \begin{center}
    {\large\bfseries Manuscrits de mathématiciens (Gallica) ---
      \ifx\grothShelfmark\empty\grothFolder\else\grothShelfmark\fi}\\[2pt]
    {\itshape\grothFoldertitle}\\[4pt]
    {\small vues \grothFirst--\grothLast
      \ifx\grothBatch\empty\else\ (lot \grothBatch)\fi}\\[2pt]
    \ifx\grothDating\empty\else
      {\small\color{gray!60!black}Datation du catalogue : \grothDating}\\[2pt]
    \fi
    \ifx\grothWatermark\empty\else
      {\small\bfseries\color{red!45!gray}\grothWatermark}\\[2pt]
    \fi
    {\footnotesize\color{gray!60!black}Transcription établie d'après les images IIIF de Gallica.
     Source gallica.bnf.fr / Bibliothèque nationale de France.\\
     Les lectures incertaines sont soulignées, les illisibles notées [\,\ldots\,],
     les ajouts entre crochets ; \textsf{v.} numérote les vues de Gallica,
     \textsf{f.} la foliotation au crayon de la BnF.}
  \end{center}
  \vspace{1em}}

\endinput


\folder{fr-9115}
\shelfmark{Français 9115}
\batch{6}
\pages{101}{120}
\dating{1801-1900}
\watermark{Lecture automatique\\non vérifiée}
\foldertitle{Papiers de Sophie GERMAIN. — Recueil de dissertations et problèmes mathématiques et physiques.}

\begin{document}

\note{Numérotation des feuillets. La série à l'encre des lots précédents se
poursuit : 50 (vue 101), 51 (vue 103), 52 (vue 105), 53 (vue 107), 54 (vue
109), 55 (vue 111), « 55 bis » au-dessus d'un « 66 » biffé (vue 113), puis,
sur les pages imprimées collées, 56 (vue 115), 57 (vue 117) et 58 (vue 119).
Les feuillets numérotés 55 et 55 bis sont écrits sur l'autre face, vues 112
et 114. Les brouillons portent en outre les lettres « A » (vue 101), « B »
(vue 103) et « D » (vue 107). L'imprimé est paginé 82 à 86 aux vues 116 à
120. Tous ces nombres ont été lus sur la vue ; aucune foliotation au crayon
de la Bibliothèque n'a été lue, et aucun « f. » n'est donc porté ici. Les
vues 102, 104, 106, 108 et 110 sont blanches (l'écriture du recto s'y voit
par transparence), les vues 111 et 113 ne portent que leur numéro : elles ne
sont pas transcrites.}

\page{101}
Le solide doué des forces d'élasticité est un système de points
matériels mobiles
\note{cette ligne de tête est écrite au-dessus du premier alinéa ; le mot
« élasticité » est repris à la plume. Elle est suivie, après « est »,
d'une lettre « A » isolée, peut-être un renvoi (la vue 95 porte
« Commencement de la p. A »). Dans la marge de gauche, en regard,
« N° 1 ». Le nombre 50 est à l'encre en haut à droite.}

Lorsque, par l'effet d'une cause extérieure le solide a changé de forme,
les points matériels qui le composent ont été déplacés : ce déplacement
ne consiste pas uniquement dans le changement de leur distance
mutuelle ; mais il se compose encore de l'espace qu'un point donné et
ceux qui l'environnent ont parcourus d'un commun mouvement.
\note{« le » est écrit par-dessus « un » au début de la deuxième ligne.}

Le solide abandonné à lui même tend à se rétablir dans sa forme
naturelle ; chacun des points mobiles qui a changé de place, est alors
animé par des forces proportionnelles \struck{\ill{}} à son déplacement
absolu et les espaces \struck{\ill{}} parcourus doivent être mesurés par
rapport à des points fixes.
\note{« des » est écrit au-dessus de « une », que la correction
remplace. Sous « proportionnelles » et sous « les espaces », deux lignes
ajoutées entre les lignes sont biffées d'un trait épais et ne se lisent
pas ; la seconde est en outre entourée au crayon, et un petit mot
« une » s'y rattache à gauche.}

Si à raison de la \struck{\uncertain{liaison}} \uncertain{symétrie} qui
\struck{unit} entre les positions des particules qui composent l'élément
du solide, \marginal{indépendamment} les espaces parcourus, dans un
même tems, par les divers points matériels, \struck{peuvent être
considérés comme un \uncertain{tout} fait et qu'on ait réussi à}
permettent de les d'exprimer, alors, la quantité et la
\uncertain{liaison} de tous \struck{ce déplacement collectif} ces
déplacemens : \struck{on connoîtra par conséquent} il sera facile
\uncertain{d'imprimer} aussi quelles sont les \marginal{simultanées}
forces qui agissent sur l'élément au moment où le solide abandonné à lui
même tend à reprendre sa forme naturelle.
\note{passage repris en tous sens, dont l'état final n'est pas cohérent.
« symétrie » est écrit au-dessus du mot biffé, « entre » au-dessus de
« unit » ; « permettent de les » au-dessus de « peuvent être
considérés », le verbe qui suivait restant sous la rature ; « alors, la
quantité et la liaison de tous » au-dessus de « ce déplacement
collectif », et « ces déplacem » dans la marge de gauche, au début de la
ligne suivante ; « il sera facile d'imprimer » sous la ligne, en regard
de « on connoîtra par conséquent » biffé. Le « à » devant « d'exprimer »
est surchargé en « d' ». Le sens attendroit « d'exprimer » dans « il sera
facile d'imprimer », mais le point sur l'\emph{i} se lit sur la vue.
« indépendamment » et « simultanées » sont écrits dans la marge de
gauche, en regard des lignes où ils sont ici placés, sans signe de
renvoi.}

\page{103}
La question des surfaces \struck{n'a} n'avoit pas été d'abord complettement
éclaircie, les expériences de Mr Chladni étoient uniquement relatives aux
cas des simples plaques. Les programmes publiés par l'accademie ont exigé
l'explication de ces expériences, le mot surface y est employé mais il
n'y est pas défini et le sens général porte à croire, que dans la pensée
des \struck{auteurs} les commissaires, la question étoit bornée à l'ordre
des faits dont la connoissance récente \struck{avoit alors frappé}
excitoit alors l'attention publique.
\note{en haut du feuillet, le nombre 51 à l'encre et, à sa droite, la
lettre « B » d'un trait appuyé — la vue 95 renvoie à « Commencement de la
p. A » et ouvre un alinéa par « À B. ». « n'avoit » est écrit au-dessus de
« n'a » biffé, « les commissaires » au-dessus de « auteurs », « excitoit »
et « alors » au-dessus des mots biffés de la dernière ligne.}

Dans le mémoire publié en 1814 par un des membres de l'accademie il ne
s'agit aussi que des seules surfaces planes.

C'est dans le mémoire couronné qui fut présenté en 1815 qu'on trouve les
premiers essais sur la théorie des surfaces courbes ; à raison de sa plus
grande simplicité on avoit \struck{choisi} traité le cas des surfaces
cylindriques. \struck{En 18\uncertain{21} j'ai publié un mémoire}

Dans mes recherches sur la théorie des surfaces publiées en 1821
\struck{je suis revenue sur} j'ai encore choisi le même exemple
j'ai cherché à généraliser l'hypothèse que j'avois proposé à l'occasion
des surfaces planes. Guidée par l'analogie j'ai \struck{clairement}
exprimé la condition relative à la direction du mouvement condition qui
renferme le caractère propre de la question.
\note{« j'ai encore choisi », suivi d'une croix de renvoi, est écrit
au-dessus de « je suis revenue sur » biffé. « j'ai cherché à généraliser
l'hypothèse » et, sous cette ligne, « que j'avois proposé à l'occasion des
surfaces planes » sont ajoutés entre les lignes, cette dernière addition
entourée d'un trait qui la rattache à « l'hypothèse » ; ils sont ici mis à
la suite. Un signe semblable à un « 3 » termine la ligne après
« l'hypothèse ». Le trait qui court sous « 1821 » et sous « le même
exemple » appartient à cet encadrement et n'est pas lu comme une rature.}

J'avoue pourtant que je ne comprenois pas alors cette question aussi
nettement que je la conçois aujourd'hui et je \uncertain{reconnoîtrai} avec
une véritable satisfaction combien les \struck{\ill{} variété} des
expériences \struck{relatives} d'autre direction de mouvemens de
Mr Savart \struck{\ill{}} par leur \uncertain{opposition} m'a aidé à
développer les conséquences des principes que j'avois déjà fixé.
\note{« relatives » (biffé, d'une encre plus pâle), « d'autre direction de
mouvemens » et « par leur opposition » sont écrits au-dessus de la ligne
qui porte « de Mr Savart » ; le mot biffé après « Savart » est effacé par
un trait épais.}

\page{105}
La discussion établie \struck{entre deux auteurs qui s'occupent de} qui
s'est à l'occasion de ce genre de questions me donne lieu de penser
\struck{que} qu'encore aujourd'hui le caractère de la surface \struck{n'a pas
été remarqué ;} n'est pas généralement \struck{elle m'engage} connue.
\struck{si on y avoit réfléchi l'objection qu'il ne s'agissoit que d'une
surface dont l'épaisseur est infiniment petite n'auroit pas été proposée
à l'auteur du mémoire publié en 1814 et ce savant géomètre ne l'auroit
pas accepté.} Je crois utile d'\ill{} \struck{sur ce que j'ai dit
ailleurs.} d'insister sur ce qu'elle Je vais essayer d'en donner une idée
complette.
\note{le nombre 52 est à l'encre en haut à droite. Le haut du feuillet
porte deux états superposés, qui ne se laissent pas démêler sûrement.
L'état premier, écrit sur les lignes et biffé ligne à ligne, puis barré
d'une grande croix sur ses trois dernières lignes, est donné d'abord ;
entre les lignes, plus serrés, viennent « qui s'est à l'occasion de »
au-dessus de « entre deux auteurs… », « qu'encore aujourd'hui » au-dessus
de « que le caractère », « n'est pas généralement ; elle m'engage » (ces
deux derniers mots biffés) au-dessus de « n'a pas été remarqué », puis
« connue. Je crois utile d’… » sous cette ligne, le dernier mot
illisible ; « sur ce que j'ai dit ailleurs. d'insister sur ce qu'elle » et
« Je vais essayer d'en donner une idée complette » dans les intervalles
suivants. L'ordre adopté ici est celui des lignes du feuillet ; « qui
s'est » n'est pas suivi du participe qu'on attendroit.}

L'expression \emph{surface} admise par une sorte de prévision longtems
avant qu'on en ait déterminé le véritable sens, désigne un solide qu'on
peut regarder comme formé par la superposition d'un nombre infini de
couches \struck{dont} qu'une l'épaisseur infiniment petite rapproche le
plus possible de l'idée mathématique de la surface.
\note{« qu'une » est écrit au-dessus de « dont » biffé ; « l'épaisseur »
n'est pas corrigé.}

Tous les points situés sur la perpendiculaire à un plan tangent donné
parcourent durant le mouvement des espaces égaux la distance entre ces
points ne change pas par l'effet du mouvement et aucun d'eux n'est
\struck{tant} porté vers les tranches voisines. Tous ces conditions qui
\uncertain{resultoit} \struck{toutes de} la direction attribuée au
mouvement qui entraîne l'invariabilité de la figure qu'affectent les
différentes couches qui composent \struck{et} l'épaisseur \struck{qui}
\struck{et de la} du solide \struck{ne seroient pas semblables} ne seroit
pas \struck{semblable} \struck{\ill{}} la même pour toute les couches et
\struck{elles} cette figure ne seroit pas non plus indépendante
\struck{\ill{}} \uncertain{ces deux fois} \struck{tout} de l'épaisseur.
Or l'expérience atteste que dans les cas de vibrations qui appartiennent
\struck{aucun genre} aux surface les figures nodales observées
\struck{sont semblables sur} sur les faces opposées sont semblables et que
l'épaisseur n'influe pas sur \struck{le déplacement des points de la
surface.} \struck{les} formation de ces figures.
\marginal{\struck{sans doute il ne faut pas confondre} Peut-être est il
inutile de remarquer que les figures qu'affectent la surface}
\note{« par l'effet du mouvement » est écrit au-dessus de « pas et aucun
d'eux » ; « qui entraîne l'invariabilité de la figure », précédé d'une
croix de renvoi et d'un « et qui entre » biffé, est porté en trois lignes
serrées au bout de la ligne « attribuée au mouvement » ; « l'épaisseur »
est écrit une première fois au-dessus de « qui composent » ; « cette
figure » au-dessus de « elles » biffé, « non plus » au-dessus de « pas
indépendante », et, sous ce mot, « ces deux fois » suivi d'un mot biffé,
dont la place dans la phrase n'est pas claire. « Or » est ajouté
au-dessus de « L'expérience », « observées » au-dessus de « sont
semblables sur », « les formation de ces figures » (le premier mot biffé) au-dessus de « le
déplacement » et sous la ligne. Le début de la ligne « l'épaisseur qui »
est couvert de taches d'encre. Les sept lignes de la note marginale sont
écrites dans la marge de gauche, en regard des cinq dernières lignes du
feuillet ; les trois premières sont biffées.}

\page{107}
pour la première fois et \struck{en} effet mon hypothèse étoit tout à fait
étrangère à la considération de l'élément du solide \struck{elle}
exprimoit uniquement la courbure qu'une cause extérieure avoit imprimée
à l'élément de la surface plane ; cependant une simple opération de
calcul a suffi pour \uncertain{en tirer} une équation \struck{qui} dont
personne ne conteste l'exactitude.
\note{en haut à droite, le nombre 53 à l'encre ; à la fin de la première
ligne, après « hypothèse », la lettre « D » d'un trait appuyé, comme le
« B » de la vue 103. Le feuillet commence au milieu d'une phrase.}

Si \struck{dans le cas des surfaces, à} par rapport aux surfaces dont la
\struck{simple} plaque \struck{est un} offre un cas particulier, le solide
ne pouvoit être légitimement regardé comme formé par la superposition
d'un nombre infini de couches d'une épaisseur infiniment petite,
assujetties à suivre dans leur mouvement toutes les circonstances du
mouvement de la couche superficielle. Comment, sans qu'il ait
\struck{rien} été statué \struck{par rapport} sur la dépendance
\struck{du mouvement} entre les déplacement des points qui composent
\struck{cette} ces couches et celui des autres points du solide
l'équation formée par la considération de la surface seule, seroit elle
devenue, à l'aide d'une simple multiplication par un coefficient
dépendant de l'épaisseur, l'expression \struck{du mouvement} des
changemens survenus dans l'élément d'un solide.
\note{« sur » est écrit au-dessus de « par rapport » biffé. Sous ce
paragraphe, un long trait horizontal à la plume.}

\struck{J'ai beaucoup insisté sur ce point parce qu'il est fondamental}
Cette notion est fondamentale en l'adoptant nous nous bornerons à
exprimer le déplacement des points qui appartiennent à la surface et
nous arriverons facilement \struck{à} ensuite à l'équation générale du
solide

\page{109}
terminé par deux faces parallèles d'une courbure quelconque sous la
réserve expresse que les mouvemens de tous les points d'un tel solide
s'exécutent dans des directions perpendiculaires aux plans tangens à la
surface qui le recouvre.
\note{le nombre 54 est à l'encre en haut à droite. Le feuillet commence au
milieu d'une phrase.}

La question relative à l'épaisseur sera traitée à part On verra plus tard
comment le cas d'une épaisseur qui ne seroit pas la même alors dans
\struck{toutes} les tranches du solide, pourvu \struck{que le mouvement ne
la fasse pas varier} toutefois qu'elle demeura invariable par l'effet du
mouvement, peut être ramené au cas d'une ou à l'égalité l'épaisseur
\struck{uniforme}. \struck{j'ai traité cette question dans un \ill{}
mémoire présenté depuis longtems à l'accademie dont on a \ill{}}
\note{« plus tard » est écrit au-dessus de « verra comment », « alors »
au-dessus de « pas la », « toutefois qu'elle demeura invariable par » et
« l'effet du mouvement » au-dessus et au-dessous de « que le mouvement ne
la fasse », « un » puis « ou à l'égalité » au-dessus de « au cas d'une »,
un petit « et » au-dessus de « uniforme » ; « j'ai traité », écrit une
seconde fois au-dessus de la ligne, est biffé comme la première. L'état
final de la phrase n'est pas arrêté sur le feuillet. Les deux dernières
lignes sont biffées d'un trait épais.}

quant à présent nous \struck{admettrons} supposerons \struck{l'épaisseur}
l'égalité \struck{cette uniformité} d'\struck{épaisseur} cette égalité
\note{« supposerons » et « l'égalité » sont écrits au-dessus de la ligne,
« d'épaisseur » (biffé) et « cette égalité » au-dessous.}

\struck{et} Après avoir établi \struck{que} qu'à l'égard des surfaces les
mouvement des points du solide seront déterminé lorsqu'on connoîtra ceux
des points qui appartiennent à la couche superficielle de ce solide
\struck{celui de la surface nous allons chercherons qu'elles sont les
espaces parcourus par les points qui environnent un point donné de
tangence.} il nous reste à montrer qu'à raison de la direction attribuée
au mouvement que les espaces parcourus \struck{durant le mouvement}, par
les points qui environnent un point donné de la surface peuvent être
commodément mesurés par le changement \struck{\ill{}}, des distances
\struck{être} \struck{survenues} les mêmes points et le plan tangent au
point donné.
\note{« qu'à l'égard » est écrit au-dessus de « que », « des surfaces »
au-dessus de « déterminé » ; « ceux des points qui appartiennent à la
couche superficielle de ce solide » occupe l'intervalle au-dessous de la
ligne ; « qui le recouvre », le deuxième mot biffé, est écrit au-dessus de
« les espaces parcourus » dans le passage biffé. « qu'à raison de la
direction » est ajouté au bout de la ligne et « attribuée au mouvement »
au-dessous, en regard de « que les espaces », sans que « que » soit biffé.
Le mot biffé dans la marge, devant « des distances », ne se lit pas ;
« survenues », biffé, est écrit au-dessus de « être ».}

La considération de la courbure moyenne qui pour le dire en passant
seroit applicable à tous les cas où la courbure \struck{est} un des
élémens admis dans les formules \struck{et ici} est ici \struck{d'une très
grande importance} de grand secours.
\marginal{\uncertain{les N°s suivans} seront \uncertain{consacrés} à
démontrer son existence et à développer ses propriétés}
\note{« en » est ajouté en tête de la ligne « passant ». La note
marginale est écrite dans la marge de gauche, en six lignes, en regard de
ce dernier alinéa.}

\page{112}
Le cas dont je me suis spécialement occupée \struck{Nous nous occuperons
uniquement du cas où tous} offre ces conditions \struck{on}
\uncertain{supposé} que les points d'un solide, dont l'épaisseur est fort
petite et dont les faces présentent une courbure quelconque, se meuvent
dans des directions perpendiculaires aux divers \struck{points} plans
tangens. Le \ill{} choix d'un de ces plans il résulte de la direction
attribuée au mouvement que tous les points situés sur la perpendiculaire
prolongée jusqu'à la face opposée du solide, parcourent, d'un commun
mouvement, des espaces égaux ; \struck{en sorte} et que si l'épaisseur
par conséquent étoit divisée en un nombre quelconque de couches
parallèles aux faces \struck{du solide} la courbure de chacune d'elles
seroit semblable à celle des faces mêmes
\note{feuillet écrit sur la face qui ne porte pas de numéro ; l'autre
face, vue 111, ne porte que le nombre 55 à l'encre, et l'écriture de
celle-ci la traverse. « Le cas dont je me suis spécialement occupée » est
écrit au-dessus de la ligne biffée, « offre ces conditions on supposé
que » au-dessous. « plans » est écrit au crayon au-dessus de « points »
biffé ; « étoit » est d'une encre plus pâle, dans la marge, et « par
conséquent » est au-dessus de « si l'épaisseur ». Le mot devant « choix »
ne se lit pas. Après « mêmes », un signe de renvoi en forme de boucle
croisée, dont l'appel n'a pas été trouvé sur le feuillet.}

Si on se borne à considérer les simples mouvemens qui \struck{d'une
amplitude fort qui s'exécutent dans} les vibrations, le cas \struck{dont}
il s'agit \struck{\ill{}} avec celui des surfaces élastiques vibrantes
\struck{et que ce n'est que cette dénomination} \struck{\ill{} que je}
\struck{et cette dénomination est d'autant plus convenable qu'elle semble
destinée à rappeler que l'épaisseur peut alors être regardé comme formé
d'une infinité de} \struck{par la superposition} \struck{\ill{}}
\struck{surface d'une épaisseur infiniment petite}
\note{« simples » est écrit au-dessus de « mouvemens », « il s'agit »
au-dessus de « dont » biffé et du mot biffé qui le suit, « élastiques
vibrantes » au-dessus de « et que ce n'est que cette ». Tout le reste de
l'alinéa, six lignes et leurs additions interlinéaires, est biffé ligne à
ligne ; au-dessus de « surface », un mot biffé ne se lit pas. Dans la
marge de gauche, en regard de « plus convenable », « ont montré » est
écrit d'une encre pâle ou au crayon, sans renvoi.}

L'identité des figures nodales qu'on observe sur les faces opposées est
une suite nécessaire de l'identité de leur courbure \struck{et}
\uncertain{confirmée} \struck{\ill{}} \struck{ingénieuses} les ingénieuses
expériences de Mr Savart \struck{lui en donne lieu de remarquer qu'elle
est}
\marginal{qu'elle est en effet le \struck{caractère} signe distinctif de
ce \struck{cas} genre \struck{cas} de vibrations}
\note{« confirmée » est écrit d'une encre pâle entre les lignes, au-dessus
de « et » ; « ingénieuses » est biffé une première fois entre les lignes.
La dernière ligne est biffée d'un trait qui la traverse, et sa lecture est
incertaine ; la note marginale, écrite en six lignes dans la marge de
gauche en regard des cinq dernières lignes, avec « en effet » d'une encre
pâle, prend la suite de la phrase. Le bas du feuillet est blanc.}

\page{114}
\[
\frac{1}{R} - \frac{1}{R} = \left(\frac{1}{r} - \frac{1}{r'}\right) \sin a \sin (a + 2A) \qquad (1)
\]
\[
\frac{1}{R} - \frac{1}{R'} = \left(\frac{1}{r} - \frac{1}{r'}\right) \sin a \sin (2A - a) \qquad (2)
\]
\[
\frac{1}{R} + \frac{1}{R''} - \frac{2}{R'} = \left(\frac{1}{r} \quad \frac{1}{r'}\right) \sin^2 a \; 2 \cos 2A \qquad (3)
\]
\note{feuillet écrit sur la face qui ne porte pas de numéro, d'une plume
plus fine et d'une encre plus pâle que les feuillets précédents, sauf les
trois équations du haut, repassées à l'encre noire et empâtées. L'autre
face, vue 113, porte en haut à droite « 55 bis » et, au-dessous, un « 66 »
biffé ; l'écriture de celle-ci la traverse. Les accents des $R$ ne sont
pas sûrs : l'encre a bavé sur le second dénominateur de l'équation (1),
dont on ne lit que le $R$, et
le signe qui sépare $\frac{1}{r}$ de $\frac{1}{r'}$ dans l'équation (3)
est couvert d'une tache, et laissé en blanc ici. Le second facteur de
l'équation (1) est écrit par-dessus une première rédaction ; celui de
l'équation (2) est surchargé et se lit « $\sin (2A - a)$ » avec doute.}

si $\frac{1}{r} = \frac{1}{r'}$, c'est à dire si la surface est une
sphère, si $a = 0$ ou $180$. si $A = 45$ c'est à dire si deux sections
font le même angle à droite et à gauche du plan de moyenne courbure avec
ce plan le second membre de l'équation (3) est nul. Nous concluerons de
là que la somme des rayon de courbure des deux sections est égale à la
somme des rayons de courbure principale
\note{« c'est à dire si la surface est une sphère » est écrit au-dessus
de la ligne ; « $A$ » est écrit sur une lettre reprise, et le « 1 » de
« 180 » sur un autre chiffre. « Nous concluerons de là que la somme des
rayon de courbure des deux sections est égale à la somme des rayons de
courbure principale » est bien ce que porte le feuillet : ce que
l'équation (3) établit porte sur les courbures, non sur les rayons.}

\struck{Se rayon dont la section dont dont le rayon est $R$ est plus près
de la section de}
\note{ces deux lignes sont barrées d'un grand trait oblique ; « rayon »
et « la section dont » y sont en outre biffés d'un trait. En regard, en
marge à droite, deux petits traits.}

D'après l'hypothèse les sections $R'$, $R$ et $R$ font les angle $A - a$,
$A$ et $A + a$ avec la section de plus grande courbure, en comparant les
seconds membres des \struck{deux} équations (1) et (2) on voit que la
différence $\frac{1}{R} - \frac{1}{R'}$ et par conséquent la diminution de
courbure entre les section $(A - a)$ et $A$ est plus petite que celle qui
a lieu entre les section $A$ et $a + A$ également distantes entr'elles que
le sont les premières. c'est à dire que \struck{la courbure} si on part de
$45^\circ$, l'augmentation de la courbure sera d'autant plus rapide pour
un même angle entre deux plans consécutifs, que le système de ces deux
plans fera un angle plus petit avec celui de plus grande courbure.
\note{« les sections $R'$, $R$ et $R$ » : les accents du second et du
troisième $R$ ne se lisent pas ; « les » est écrit sur un autre mot. Le
« $A + a$ » de la première ligne est écrit sur un signe repris.}

Les deux parties de la courbe qui marque les différences De courbure
entre les différentes sections normales d'un même élément sont,
symétriques autour de l'ordonnée de moyenne courbure une de ces parties,
celle entre la section de plus grande courbure et la section de $45^\circ$,
est convexe du côté de \struck{l'a} diamètre qui tient la place du plan
tangent l'autre est concave vers le même diamètre. La courbure diminuera
donc d'autant plus pour un même angle entre deux sections consécutives que
l'on s'approchera davantage de \struck{l'ang} la section de moyenne
courbure
\note{« qui » est écrit sur une lettre reprise ; « la » devant
« section », à la quatrième ligne, est d'une encre plus pâle. Le reste du
feuillet est blanc.}

\section{Pièce imprimée : mémoire de Legendre sur la détermination des fonctions Y et Z}

\page{115}
\note{feuillet imprimé, collé sur un feuillet de montage : première page du
mémoire. Le nombre 56 est à l'encre en haut à droite du feuillet ; un
timbre rond de la « Bibliothèque impériale » est apposé à côté du nom de
l'auteur. La
page ne porte pas de numéro imprimé ; au pied, « T. XI. » à gauche et la
signature « 11 » au milieu.}

MÉMOIRE

SUR

\emph{La détermination des fonctions Y et Z qui satisfont à l'équation
$4(x^n - 1) = (x - 1)(Y^2 \pm nZ^2)$, n étant un nombre premier
$4i \mp 1$.}

PAR M. LEGENDRE.

Lu à l'Académie, le 11 octobre 1830.

Dans les articles 511 et 512 de ma Théorie des nombres, troisième édition,
j'ai donné des règles fort simples pour calculer les coefficients de la
fonction Y qui satisfait à l'équation $4X = Y^2 \pm nZ^2$, où l'on désigne
par $n$ le nombre premier $4i \mp 1$ et par X le quotient de $x^n - 1$
divisé par $x - 1$. Ces règles sont fondées sur la supposition que les
coefficients des différents termes de la fonction Y sont moindres que
$\frac{1}{2}n$, en les prenant avec le signe convenable ; cette propriété a
lieu en effet pour toutes les valeurs de $n$ comprises dans le tableau de
l'art. 512, et elle se vérifierait également pour quelques-unes des valeurs
suivantes, mais elle cesse d'avoir lieu pour des valeurs plus grandes et il
arrive au contraire que les coefficients de $x$ dans la fonction Y, loin
d'être plus

\page{116}
\note{p. 82 imprimée, verso du feuillet précédent, au titre courant
« MÉMOIRE ». Le signe à l'encre en haut à gauche est le « 56 » de la vue
115 vu par transparence.}

petits que $\frac{1}{2}n$, deviennent progressivement plus grands que $n$,
à mesure que $n$ augmente, même plus grands que $n^2$, $n^3$, etc. Les
règles dont nous parlons ne font donc plus connaître les vrais coefficients
de Y, mais seulement les restes de ces coefficients divisés par $n$, ce qui
peut être utile pour vérifier les valeurs calculées par d'autres méthodes,
et il en résulte toujours, conformément à ce qui a été dit dans l'article
511, que le polynome Y, dont les premiers termes sont
$2x^m + x^{m-1} + \frac{3 \mp n}{4}x^{m-2} + \text{etc.}$ contient
nécessairement tous ses termes au nombre de $m + 1$ sans qu'aucun d'eux
puisse devenir nul, ce qui n'a pas toujours lieu pour le polynome Z du
degré $m - 1$, qui peut être complet ou incomplet suivant les différents
cas. Voici maintenant une méthode sûre et exacte pour déterminer, dans tous
les cas, les polynomes Y et Z qui satisfont à l'équation
$4X = Y^2 \pm nZ^2$.

Soit d'abord $n = 4i + 1$ ou $m = 2i$, valeur qui répond à l'équation
$4X = Y^2 - nZ^2$, si on fait en général
\[
Y + Z\sqrt{n} = 2x^m + A_1 x^{m-1} + A_2 x^{m-2} + A_3 x^{m-3} + \text{etc.}
\]
la forme connue des polynomes Y et Z étant
\[
\begin{array}{l}
Y = 2x^m + x^{m-1} + a_2 x^{m-2} + a_3 x^{m-3} + \text{etc.} \\
Z = \phantom{2x^m + {}} x^{m-1} + b_2 x^{m-2} + b_3 x^{m-3} + \text{etc.},
\end{array}
\]
on aura les relations suivantes entre les coefficients $A$, $a$, $b$ :
\[
A_1 = 1 + \sqrt{n}, \quad A_2 = a_2 + b_2\sqrt{n}, \quad A_3 = a_3 + b_3\sqrt{n},
\]
et en général $A_k = a_k + b_k\sqrt{n}$ ; ainsi pour avoir les valeurs des
fonctions Y et Z, il suffira de connaître celles des coefficients $A_2$,
$A_3$, etc.

\page{117}
\note{p. 83 imprimée, au titre courant « SUR UNE QUESTION D'ANALYSE. », sur
un second feuillet collé sur son propre feuillet de montage. Le nombre 57
est à l'encre en haut à droite ; au pied, la signature « 11. ».}

Soit $g$ l'une des racines primitives du nombre premier $n$ : si l'on
désigne par $\alpha$ l'un des termes de la suite
\[
1,\ g^2,\ g^4,\ g^6 \ldots\ldots g^{n-3}, \qquad (1)
\]
et par $\beta$ l'un des termes de la suite
\[
g,\ g^3,\ g^5,\ g^7 \ldots\ldots g^{n-2},
\]
on sait que ces deux suites de nombres, diminués des multiples de $n$
qu'ils peuvent contenir, donnent par leur réunion la suite des nombres
naturels $1, 2, 3 \ldots\ldots n - 1$ ; on sait de plus que $r$ étant
l'une quelconque des racines imaginaires de l'équation $r^n - 1 = 0$, le
produit de tous les facteurs $x - r^\alpha$, sera égal à l'un des deux
polynomes $\frac{1}{2}(Y + Z\sqrt{n})$, $\frac{1}{2}(Y - Z\sqrt{n})$, et le
produit de tous les facteurs $x - r^\beta$ sera égal à l'autre polynome ;
d'ailleurs comme le signe de $\sqrt{n}$ peut être pris à volonté, on
pourra supposer généralement
\[
\begin{array}{l}
(x - r)(x - r^{g^2})(x - r^{g^4}) \ldots (x - r^{g^{n-3}}) \\
\quad = x^m + \frac{1}{2}A_1 x^{m-1} + \frac{1}{2}A_2 x^{m-2} + \frac{1}{2}A_3 x^{m-3} + \text{etc.}
\end{array}
\]

Appelons $S_1$ la somme des racines $r^\alpha$, $S_2$ la somme de leurs
carrés $r^{2\alpha}$, $S_3$ la somme de leurs cubes $r^{3\alpha}$, etc. ; ces
sommes étant supposées connues, on pourra déterminer les coefficients
$A_1$, $A_2$, $A_3$, etc., au moyen des équations suivantes :

(1) La suite des valeurs de $\alpha$ est aussi représentée dans un autre
ordre, par celle des carrés $1, 4, 9, 16 \ldots m^2$, diminués des multiples
de $n$ qu'ils peuvent contenir.
\note{note infrapaginale de l'imprimé, sous un filet ; son appel est le
« (1) » placé à droite de la première suite.}

\page{118}
\note{p. 84 imprimée, au titre courant « MÉMOIRE », verso du feuillet
précédent ; le signe à l'encre en haut à gauche est le « 57 » de la vue
117 vu par transparence.}

\[
(1) \qquad
\begin{array}{l}
-\frac{1}{2}A_1 = S_1 \\
-2.\frac{1}{2}A_2 = S_2 + S_1.\frac{1}{2}A_1 \\
-3.\frac{1}{2}A_3 = S_3 + S_2.\frac{1}{2}A_1 + S_1.\frac{1}{2}A_2 \\
-4.\frac{1}{2}A_4 = S_4 + S_3.\frac{1}{2}A_1 + S_2.\frac{1}{2}A_2 + S_1.\frac{1}{2}A_3 \\
-5.\frac{1}{2}A_5 = S_5 + S_4.\frac{1}{2}A_1 + S_3.\frac{1}{2}A_2 + S_2.\frac{1}{2}A_3 + S_1.\frac{1}{2}A_4 \\
\text{etc.}
\end{array}
\]

A l'égard de la somme désignée en général par $S_k$, elle ne peut avoir
que l'une des deux valeurs $-\frac{1}{2} + \frac{1}{2}\sqrt{n}$,
$-\frac{1}{2} - \frac{1}{2}\sqrt{n}$.

On peut supposer $S_1 = -\frac{1}{2} - \frac{1}{2}\sqrt{n}$, et alors on
aura
\[
S_k = -\tfrac{1}{2} - \tfrac{1}{2}\sqrt{n}, \text{ si } k \text{ appartient à la série des nombres } \alpha ;
\]
et
\[
S_k = -\tfrac{1}{2} + \tfrac{1}{2}\sqrt{n}, \text{ si } k \text{ appartient à la série des nombres } \beta .
\]

Dans le premier cas, le nombre $k$ serait un résidu quadratique de $n$, et
on aurait, suivant la notation ordinaire, $\left(\frac{k}{n}\right) = 1$ ;
dans le second cas, le nombre $k$ serait un non-résidu et on aurait
$\left(\frac{k}{n}\right) = -1$ ; donc dans tous les cas on aura, au moyen
du symbole $\left(\frac{k}{n}\right)$,
\[
(2) \qquad S_k = -\tfrac{1}{2} - \tfrac{1}{2}\sqrt{n}.\left(\frac{k}{n}\right).
\]

Étant donnés $n$ et $k$, on pourra toujours déterminer \emph{a priori}
celle des deux valeurs $+1$ et $-1$ qui convient au symbole
$\left(\frac{k}{n}\right)$. Ainsi on connaîtra successivement toutes les
sommes $S_1$, $S_2$, $S_3$, etc., ce qui permettra de déterminer les
coefficients $A_1$, $A_2$, $A_3$, par le moyen des équations (1). Ensuite
la fonction $2x^m + A_1 x^{m-1} + A_2 x^{m-2} + A_3 x^{m-3} + \text{etc.}$
prendra la forme
\note{le bas de la page imprimée est blanc : le texte reprend à la p. 85.}

\page{119}
\note{p. 85 imprimée, au titre courant « SUR UNE QUESTION D'ANALYSE. », sur
un troisième feuillet collé sur son feuillet de montage. Le nombre 58 est à
l'encre en haut à droite.}

\[
\begin{array}{l}
2x^m + (1 + \sqrt{n})x^{m-1} + (a_2 + b_2\sqrt{n})x^{m-2} \\
\qquad + (a_3 + b_3\sqrt{n})x^{m-3} + \text{etc.}
\end{array}
\]
d'où l'on tire les valeurs séparées des fonctions Y et Z, savoir :
\[
\begin{array}{l}
Y = 2x^m + x^{m-1} + a_2 x^{m-2} + a_3 x^{m-3} + \text{etc.} \\
Z = \phantom{2x^m + {}} x^{m-1} + b_2 x^{m-2} + b_3 x^{m-3} + \text{etc.}
\end{array}
\]

Telle est la méthode par laquelle on trouvera les fonctions Y et Z qui
satisfont à l'équation $4X = Y^2 - nZ^2$, lorsque le nombre premier $n$ est
de la forme $4i + 1$ ; elle servira également à résoudre l'équation
$4X = Y^2 + nZ^2$, lorsque le nombre premier $n$ est de la forme $4i - 1$.
Il suffira pour cela de mettre $-n$ à la place de $n$ dans la valeur de
$S_1$, c'est-à-dire, de prendre
$S_1 = -\frac{1}{2} - \frac{1}{2}\sqrt{-n}$ ; on déterminera ensuite
généralement $S_k$ au moyen de la formule
$S_k = -\frac{1}{2} - \frac{1}{2}\sqrt{-n}.\left(\frac{k}{n}\right)$ en y
substituant la valeur particulière de $\left(\frac{k}{n}\right)$ qu'on peut
toujours trouver \emph{a priori}. Alors connaissant par les équations (1)
les valeurs des coefficients $A$, $A_2$, $A_3$, etc. qui seront tous de la
forme $a + b\sqrt{-n}$, on en déduira comme ci-dessus la valeur de la
fonction $2x^m + A_1 x^{m-1} + A_2 x^{m-2} + \text{etc.}$, et ensuite celles
des fonctions Y et Z.
\note{« $A$, $A_2$, $A_3$ » : l'indice du premier $A$ manque dans
l'imprimé.}

\subsection{Exemple I.}

Pour continuer le tableau de l'art. 512, prenons $n = 31$. Comme les
diverses valeurs de $S_k$ ne peuvent être que P et Q, savoir
$P = -\frac{1}{2} - \frac{1}{2}\sqrt{n}$,
$Q = -\frac{1}{2} + \frac{1}{2}\sqrt{n}$, ayant déja fait $S_1 = P$, nous
déterminerons les sommes suivantes $S_2$, $S_3$, $S_4$, etc. au moyen des
symboles $\left(\frac{2}{n}\right)$, $\left(\frac{3}{n}\right)$,
$\left(\frac{4}{n}\right)$, $\left(\frac{5}{n}\right)$, etc.
\note{le bas de la page imprimée est blanc. Pour $n = 31$, de la forme
$4i - 1$, la règle qui précède demanderoit $\sqrt{-n}$ ; l'imprimé écrit ici
$\sqrt{n}$, et c'est ce qui est transcrit.}

\page{120}
\note{p. 86 imprimée, au titre courant « MÉMOIRE », verso du feuillet
précédent ; le signe à l'encre en haut à gauche est le « 58 » de la vue
119 vu par transparence.}

Or on trouve par les principes connus
\[
\begin{array}{l}
\left(\frac{2}{31}\right) = 1,\ \left(\frac{3}{31}\right) = -\left(\frac{31}{3}\right) = -1,\ \left(\frac{4}{31}\right) \text{ ou en général } \left(\frac{c^2}{31}\right) = 1, \\
\quad \left(\frac{5}{31}\right) = \left(\frac{31}{5}\right) = 1,\ \left(\frac{6}{31}\right) = \left(\frac{2}{31}\right).\left(\frac{3}{31}\right) = -1,\ \left(\frac{7}{31}\right) = -\left(\frac{31}{7}\right) \\
\quad = -\left(\frac{3}{7}\right) = \left(\frac{7}{3}\right) = 1,\ \left(\frac{8}{31}\right) = \left(\frac{2c^2}{31}\right) = \left(\frac{2}{31}\right) = 1,\ \left(\frac{9}{31}\right) = 1, \\
\quad \left(\frac{10}{31}\right) = \left(\frac{2}{31}\right).\left(\frac{5}{31}\right) = 1, \text{ etc.},
\end{array}
\]
on a donc
\[
\begin{array}{l}
S_1 = S_2 = S_4 = S_5 = S_7 = S_8 = S_9 = S_{10} = P \\
S_3 = S_6 = Q.
\end{array}
\]

Ensuite par l'application des équations (1) on trouvera
\[
\begin{array}{l}
A_1 = 1 + \sqrt{-n} \\
A_2 = -7 + \sqrt{-n} \\
A_3 = -11 - \sqrt{-n} \\
A_4 = 2 - 2\sqrt{-n} \\
A_5 = 8 \\
A_6 = -3 + \sqrt{-n} \\
A_7 = -5 - \sqrt{-n} \\
A_8 = 5 - \sqrt{-n}.
\end{array}
\]

Arrivé aux termes moyens $A_7$ et $A_8$ on peut se dispenser d'aller plus
loin, car il est facile de voir qu'on doit changer simplement le signe de la
partie réelle dans les coefficients précédents à compter de $A_7$, ce qui
donnera
\note{le bas de la page imprimée est blanc ; la suite de l'exemple est sur
la vue 121, hors de ce lot.}

\end{document}
